\documentclass[10pt,letterpaper]{article}
\usepackage{cornell-notes}

\title{Clause 07.06.20: Comma Operator}
\author{}
\date{}
\setDocTitle{Clause 07.06.20: Comma Operator}
\setDocSubtitle{C++ 2024}
\setDocOwner{C++ 2024 Cornell Notes}
\setDocDate{\today}
\setCornellCollection{C++ 2024 Cornell Notes}
\setCornellUnitType{Clause}
\setCornellUnitNumber{07.06.20}
\setCornellUnitTitle{Comma Operator}
\hypersetup{pdftitle={Clause 07.06.20: Comma Operator},pdfsubject={Cornell notes},pdfauthor={}}

\begin{document}
\maketitle
\begin{CornellOverviewBox}{Overview}
\textbf{Classification:} Core language - Normative\\[2pt]
\textbf{Learning objective:} Understand the rules, terminology, and semantic consequences associated with Comma operator.
\end{CornellOverviewBox}\begin{CornellSummaryBox}{Summary}
{\sffamily\bfseries\color{Accent} Summary}\par\vspace{3pt}Section 7.6.20 focuses on Comma operator. Use the listed grammar productions together with the semantic constraints and disambiguation rules referenced by the section. When applying it, preserve the distinction between mandatory requirements, recommendations, permissions, and behavior deliberately left undefined, unspecified, or implementation-defined.
\end{CornellSummaryBox}

\end{document}
